% !TeX root = ../thuthesis-example.tex

\chapter{后续研究计划}

当前论文版本已经完成了与代码实现一致的路径级建模、非保序性证明和 \texttt{Dijkstra} / \texttt{Sliced Dijkstra} 算法对比。后续工作将围绕实验验证与论文整合继续推进。

\section{第一阶段：模型与证明定稿}

继续核对 \texttt{BottleneckFidelityRoutingMetric}、\texttt{DijkstraRoutingProtocol} 和 \texttt{SlicedDijkstraRoutingProtocol} 的实现细节，确保论文中的符号、公式和伪代码与当前代码一致。重点检查 $T_{\max}$ 对齐、相位翻转等待退相干、Bell 对角态交换以及 \texttt{IsBetter} / \texttt{Dominates} 的排序语义。

\section{第二阶段：手工拓扑验证}

构造小规模前缀陷阱拓扑，验证如下现象：

\begin{enumerate}
  \item 单标签 \texttt{Dijkstra} 在中间节点处只保留局部分数最高的 label；
  \item 局部分数更高但 $T_{\max}$ 更大的前缀，在接入共同后缀后可能被反超；
  \item \texttt{Sliced Dijkstra} 在不同 $T_{\max}$ bucket 中保留多个前缀，从而有机会恢复被单标签剪枝的候选。
\end{enumerate}

\section{第三阶段：随机拓扑与真实拓扑实验}

在随机拓扑中系统扫描 $K$、\texttt{BucketWidthMs} 和 \texttt{UseBuckets} 对路由结果的影响，观察切片算法相对于单标签基线的路由差异率和预测保真度差异。在真实拓扑实验中，以 Topology Zoo 等网络为基础构造量子-经典混合拓扑，并对路由结果进行混态仿真复核。

需要注意的是，切片算法不是严格最优算法，因此实验分析不应只报告“平均提升”，还应记录失败样本和原因，例如有限 $K$ 截断、同桶局部代表失效、路径跳数增加导致实际交换等待变长等。

\section{第四阶段：论文整合与结果补充}

在实验完成后，将当前占位式实验设计替换为正式实验章节。结果叙述将区分三类量：

\begin{enumerate}
  \item 路由阶段的 \texttt{predicted\_fidelity}；
  \item 混态仿真得到的实际端到端保真度；
  \item 算法搜索行为本身，例如 label 保留数、bucket 分布和被剪枝前缀。
\end{enumerate}

这样可以避免把路由 metric 的预测增益和完整物理仿真的实际增益混为一谈，也能更清楚地解释切片算法在什么条件下有效、在什么条件下失效。
